\section{Supplementary Problem 2.4}
The task is to show that every discrete-time signal can be decomposed
into an even and an odd component, determine whether this decomposition
is unique, and illustrate the result with $x(n)=\{2,3,4,5,6\}$, where
the central sample $4$ corresponds to $n=0$.

\subsection{Existence of the decomposition}
For an arbitrary signal $x(n)$, define
\begin{align}
  x_e(n)&=\frac{x(n)+x(-n)}{2},\\
  x_o(n)&=\frac{x(n)-x(-n)}{2}.
\end{align}
Evaluating at $-n$ gives
\begin{align}
  x_e(-n)&=\frac{x(-n)+x(n)}{2}=x_e(n),\\
  x_o(-n)&=\frac{x(-n)-x(n)}{2}=-x_o(n).
\end{align}
Thus $x_e$ is even and $x_o$ is odd. Their sum is
\begin{equation}
  x_e(n)+x_o(n)
  =\frac{x(n)+x(-n)+x(n)-x(-n)}{2}=x(n),
\end{equation}
which proves that the decomposition exists for every signal.

\subsection{Uniqueness of the decomposition}
Suppose $x(n)=a(n)+b(n)$ for some even signal $a$ and odd signal $b$.
Their symmetry implies
\begin{equation}
  x(-n)=a(-n)+b(-n)=a(n)-b(n).
\end{equation}
Adding and subtracting the two equations forces
\begin{align}
  a(n)&=\frac{x(n)+x(-n)}{2}=x_e(n),\\
  b(n)&=\frac{x(n)-x(-n)}{2}=x_o(n).
\end{align}
Every even/odd decomposition therefore has exactly these components,
so the decomposition is unique. This algebraic argument applies to
arbitrary $x$; the numerical example below only illustrates the result.

\subsection{Example and expected result}
The given sequence is stored on $n=-2,-1,0,1,2$, with zero values
outside that interval. Since this time grid is symmetric about zero,
reversing the vector gives the samples of $x(-n)$ on the same grid.
\begin{table}[H]
  \centering
  \caption{Unique even and odd components for Problem 2.4}
  \label{tab:p24}
  \begin{tabular}{lrrrrr}
    \toprule
    Time $n$ & $-2$ & $-1$ & 0 & 1 & 2\\
    \midrule
    $x(n)$ & 2 & 3 & 4 & 5 & 6\\
    $x(-n)$ & 6 & 5 & 4 & 3 & 2\\
    $x_e(n)$ & 4 & 4 & 4 & 4 & 4\\
    $x_o(n)$ & $-2$ & $-1$ & 0 & 1 & 2\\
    $x_e(n)+x_o(n)$ & 2 & 3 & 4 & 5 & 6\\
    \bottomrule
  \end{tabular}
\end{table}
Consequently,
\begin{align}
  x_e(n)&=\begin{cases}4,& |n|\leq2,\\0,&\text{otherwise},\end{cases}\\
  x_o(n)&=\begin{cases}n,& |n|\leq2,\\0,&\text{otherwise}.\end{cases}
\end{align}
Both vectors retain origin position three. The even component is
constant only on the stated finite interval, and the odd component
has a zero sample at the origin as required.

\subsection{Scilab implementation}
The script computes the two components and their sum, then displays
four plots: the original signal, the even component, the odd component,
and the reconstructed signal. It also prints three numerical checks:
\begin{align}
  \max_n|x(n)-x_e(n)-x_o(n)|&=0,\\
  \max_n|x_e(n)-x_e(-n)|&=0,\\
  \max_n|x_o(n)+x_o(-n)|&=0.
\end{align}
These are expected results for the specified example; the proof of
uniqueness is the general derivation above.
\lstinputlisting[caption={Scilab implementation of supplementary Problem 2.4},label={lst:p24}]{../2_4.sci}

\subsection{Execution evidence}
\resultfigure{images/problem24.png}
  {Original signal, even and odd components, and reconstruction for Problem 2.4}
  {fig:p24}
The captured plots match \Cref{tab:p24}. The even component has five
samples of amplitude four, the odd component has amplitudes
$-2,-1,0,1,2$, and the reconstructed signal matches the original.
